% ═══════════════════════════════════════════════════════════════════
% ЧАСТЬ I. ПРОГРАММА И ПРИНЦИПЫ
% ═══════════════════════════════════════════════════════════════════
\part{Программа и принципы}

\section{Введение: динамический принцип}\label{sec:intro}

Настоящая монография объединяет в единое целое результаты программы,
развивавшейся поэтапно через серию вычислительных стендов: калибровочный
тор, поверхность K3, якобиан квартики Клейна и циклотомический стенд
$N=15/30$ на кривых Ферма. Каждая ступень этой лестницы была описана в
отдельных материалах; здесь все они собраны, согласованы между собой и
дополнены полными выкладками, так что монография читается как одно
связное исследование — от исходного принципа до замкнутой системы
лемм и теорем с воспроизводимым вычислительным протоколом. Все
численные утверждения снабжены сертификатами: точными целочисленными
проверками либо прямым независимым интегрированием, воспроизводимым
одной командой из лаборатории \texttt{laboratory.py}.

Основание программы составляет \emph{динамический принцип} — правило,
согласно которому вся структура поправок строится из целочисленных
геометрических данных. Принцип был зафиксирован при аудите исходных
монографий и с тех пор ни разу не нарушался; его формулировка проста.

\begin{keybox}{Динамический принцип (главная формулировка программы)}
\emph{Геометрия поставляет только целочисленную тройку
$(n, B, \lambda_0)$: порядок $n$ выделенного элемента симметрии,
целочисленный топологический индекс $B$ носителя и спектральный порог
$\lambda_0$. Вся остальная структура — фазы, торможения, эффективные
поправки, нормировки, решётки периодов — строится из этой тройки
чистой алгеброй и не содержит свободных параметров.}
\end{keybox}

Принцип силён именно своей скупостью: геометрия ничего не «угадывает»
и не подбирает коэффициентов. Она лишь отвечает на три вопроса:
какой порядок имеет поворот, каков целый индекс носителя и каков порог
спектра. Дальше работает универсальный алгебраический конвейер, один и
тот же на всех ступенях лестницы. Именно этим объясняется единообразие
формул на торе, на K3, на квартике Клейна и на кривых Ферма: формулы
не повторяются по аналогии — они выводятся из одной и той же тройки
по одним и тем же правилам.

Важным методологическим решением программы является отказ от методов,
которые могли бы создать видимость результата там, где его нет. В
окончательной редакции не используются: метод PSLQ и любые
целочисленно-реляционные подгонки для утверждений основного текста;
L-значения характеров Гекке и «общие фазы», постулируемые извне;
численные проверки того, что доказывается тождественно; апелляции к
самосопряжённым операторам, которые в конструкции не возникают. Там,
где численная сверка нужна, она выполняется независимым методом —
прямым квадратурированием по гладким кускам без $\Ga$-функций — и
служит сертификатом, а не доказательством.

\subsection{Что доказано и что вычислено}

Программа накопила замкнутую систему результатов, каждый из которых
доказан в тексте настоящей монографии. Для ориентировки читателя
приведём сводку; полные формулировки и доказательства занимают
соответствующие разделы.

\begin{statusbox}[сводка программы]
\begin{itemize}[leftmargin=1.4em, itemsep=0.15em]
\item Род кривой Ферма $g=\tfrac{(N-1)(N-2)}{2}$ и тройка тора
  $(4,1,4\pi^2)$ — доказаны из первых принципов (разделы~\ref{sec:torus},
  \ref{sec:cycles}).
\item Ранг Нерона--Севери $\rho=20$ для Фермат-квартики $x^4+y^4+z^4+w^4=0$,
  сигнатура $(1,19)$, определитель подрешётки $64$ — вычислены точно над
  $\QQ$ по явной матрице пересечений $49\times49$ (раздел~\ref{sec:k3}).
\item Errata E8: на роде $3$ из $64$ спиновых структур чётных $36$ ($\mathrm{Arf}=0$),
  нечётных $28$ ($\mathrm{Arf}=1$) — доказано полным перечислением по инварианту Арфа
  (раздел~\ref{sec:e8}).
\item Сертификаты B, C, D, E, F, G, H — приняты по всем пунктам своих
  протоколов: закрытие ядра $29$ периодами; $\mu_4$-эквивариантность
  поправок $22=1+7+7+7$; универсальная $\mu_4$-теорема $[L,P]=0$;
  завершимость потока с точным временем $t^*=\lcm(\dots)$;
  рационализация с априорной границей знаменателей $Q=256$;
  индексы решётки Ходжа и слой Коула--Селберга; стенд $N=15/30$ в полной
  $\Ga$-нормировке Гросса (разделы~\ref{sec:certB}--\ref{sec:klein}).
\item Замкнутая форма периода $P_{a,b}(r,s)=\tfrac1N\zeta_N^{ra+sb}
  \Omega_{a,b}$, соотношения Римана как тождество, $\Ga$-нормировка с
  дробными долями $\pi$, поляризация и вычислимость индексов —
  замкнутая система лемм и теорем стенда $N=15/30$
  (разделы~\ref{sec:cycles}--\ref{sec:indices}).
\item Протокол V1--V9: перепись характеров $91=1+6+84$ и
  $406=1+6+9+30+84+276$; отклонения замкнутой формы от независимого
  интегрирования $3.9\cdot10^{-36}$; фазы с отклонением $0.0$;
  ранг системы $=g$ при числах обусловленности $8.6$ и $18.2$;
  глубокая запись при $dps=70$ с отклонением $1.7\cdot10^{-71}$
  (раздел~\ref{sec:protocol}).
\end{itemize}
\end{statusbox}

\subsection{Как устроена монография}

Изложение организовано пятью частями. Часть~I (настоящий раздел) вводит
принцип, три поправки и методологию лестницы. Часть~II описывает
калибровочные стенды: тор, на котором вся механика видна в явном виде;
поверхность K3 с максимальным рангом Нерона--Севери; errata E8 по
спиновым структурам рода~3; бинарный код как мост между геометрией и
циклической алгеброй. Часть~III собирает сертификаты программы — от
сертификатора циклов (сертификат~A) до сертификата~G об индексах
решётки Ходжа; сюда же входит третий стенд — якобиан квартики Клейна.
Часть~IV посвящена стенду $N=15/30$: это самая зрелая часть программы,
где изложение доведено до замкнутой системы лемм и теорем, не
опирающейся ни на что внешнее. Часть~V содержит протокол V1--V9, сводные
таблицы и приложения с полными выводами.

Читателю, который хочет сначала увидеть механику в действии,
рекомендуется начать с раздела~\ref{sec:torus} о торе: там вывод
поправок занимает полторы страницы и каждая величина видна явно.
Читателю, интересующемуся фундаментом, следует идти подряд: комбинаторика
циклов (раздел~\ref{sec:cycles}) построена без всякой аналитики и
служит основанием всего дальнейшего.

\section{Три поправки: происхождение и обоснование}\label{sec:corrections}

Ядро фреймворка составляют три поправки, каждая из которых выводится,
а не постулируется. В этом разделе мы фиксируем их определения,
обосновываем происхождение каждой и собираем объединённую формулу.
Ни одна из величин не является подгоночным параметром: все они —
функции целочисленной тройки $(n,B,\lambda_0)$.

\subsection{Спиновая фаза $\delta=\pi/N$}

Первая поправка — спиновая фаза. Пусть геометрия поставила элемент
симметрии порядка $N$: поворот, перестанавливающий ветви спинового
двойного накрытия. В спин-расслоении подъём такого поворота имеет
порядок $2N$ и переставляет два листа накрытия; фаза, приобретаемая
собственной формой при обходе оси поворота, равна половине элементарного
угла:
\begin{equation}\label{eq:delta}
  \delta \;=\; \frac{\pi}{N}.
\end{equation}
Происхождение формулы — двойное накрытие: на базовой кривой поворот
имеет период $N$, но в спин-расслоении полный обход оси возвращает
лист лишь после двух базовых обходов, поэтому элементарный фазовый шаг
равен $\pi/N$, а не $2\pi/N$. На кривой Клейна порядок образующей
$C$ группы $\Gamma(2,3,7)$ равен $7$, и фаза $\delta_C=\pi/7$; на торе
порядок поворота равен $4$, и фаза $\delta=\pi/4$; на кривых Ферма
роль $N$ играет уровень, и фазы кратны $\pi/15$ и $\pi/30$.

\begin{remark}[почему фаза не «подбирается»]
Формула~\eqref{eq:delta} не содержит свободных множителей. Любое
отклонение от неё означало бы, что элемент симметрии имеет порядок,
отличный от $N$; но порядок — целочисленная характеристика геометрии.
Таким образом, спиновая фаза наследует целочисленность геометрии:
все дальнейшие поправки строятся из $\delta=\pi/N$ алгебраически.
\end{remark}

\subsection{Торможение $\gamma=\delta^4/k$ и эффективная фаза
$\delta_{\mathrm{eff}}=\delta^5/k$}

Вторая поправка — торможение. Поправка порядка четыре в фазе
подавляется носителем: роль делителя играет целочисленный
топологический индекс $B$ — число $b_2$ многообразия-носителя либо
другой целочисленный индекс ступени. Определение:
\begin{equation}\label{eq:gamma}
  \gamma \;=\; \frac{\delta^4}{k},
  \qquad k = B\ \text{(индекс носителя)}.
\end{equation}
Универсальный выбор фреймворка — $k=b_2(K3)=22$: второе число Бетти
поверхности K3 является целочисленной константой, общей для всех
ступеней лестницы, поскольку K3-слой присутствует в конструкции как
носитель поправок. На торе носителем служит сам тор ($k=1$), и торможение
принимает значение $\gamma=\delta^4$; на кривой Клейна используется
$k=22$; на стенде $N=15/30$ индексов несколько — по блокам проводников,
и они вычисляются, а не постулируются (раздел~\ref{sec:indices}).

Третья поправка объединяет фазу с торможением в одну эффективную
величину:
\begin{equation}\label{eq:deff}
  \delta_{\mathrm{eff}} \;=\; \frac{\delta^5}{k} \;=\; \delta\cdot\gamma.
\end{equation}
Для кривой Клейна $(\pi/7)^5/22\approx1/1200$ — поправка мала, что и
требуется: эффективная фаза должна быть возмущением, а не заменой
главного члена. На торе при $k=1$ эффективная фаза велика
($(\pi/4)^5\approx0.299$), и именно поэтому тор служит калибровочным
стендом: все эффекты видны при грубом счёте, и любые ошибки конструкции
проявляются немедленно.

\subsection{Объединённая формула $\Delta_{Ch}$}

Три поправки собираются в одну величину вместе со спектральным порогом
$\lambda_0$ и ранговой добавкой $R/4$:
\begin{keybox}{Объединённая формула}
\begin{equation}\label{eq:deltaCh}
  \Delta_{Ch} \;=\; \lambda_0 \;-\; \frac{R}{4}
  \;+\; \frac{\delta^2}{2} \;-\; \frac{\delta^5}{k},
  \qquad
  b_{Ch}(n) \;=\; 1-\cos\frac{2\pi}{n}.
\end{equation}
Здесь $\lambda_0$ — спектральный порог тройки; $R$ — целочисленный
ранговый параметр ступени; $\delta=\pi/n$ — спиновая фаза; $k$ —
индекс носителя торможения. Константа $b_{Ch}(n)$ — радикальная
характеристика поворота порядка $n$, на уровнях стенда принимающая
точные радикальные значения (теорема~\ref{th:normalization}).
\end{keybox}

Каждый член формулы~\eqref{eq:deltaCh} имеет собственное происхождение.
Член $\lambda_0$ — «голое» собственное значение, поставляемое
геометрией; для кривой Клейна это $\lambda_1(D^2_{\sigma_0})=3.338$,
для тора $4\pi^2$. Член $-R/4$ — ранговая добавка, отражающая размер
нерасщеплённой части носителя. Член $+\delta^2/2$ — главная фазовая
поправка второго порядка, универсальная для всех поворотов. Член
$-\delta^5/k$ — торможение пятого порядка: малая, но обязательная
коррекция, отличающая фреймворк от «чистого» $\delta^2/2$.

Согласование с наблюдаемыми значениями проверено на всех ступенях.
Для кривой Клейна объединённая формула даёт
$\Delta_{Ch}\approx3.4399$ при наблюдаемом значении $3.443$ —
отклонение $0.0905\%$, зафиксированное в данных протокола
(таблица~\ref{tab:torus-triples}). Отклонение не устраняется подгонкой:
оно есть следствие одних и тех же правил для всех ступеней, и его
малость служит свидетельством правильности правил, а не настройки.

\section{Методология лестницы стендов}\label{sec:ladder}

Лестница стендов — это последовательность геометрических объектов
возрастающей сложности, на каждом из которых один и тот же конвейер
запускается целиком. Лестница устроена так, что каждая следующая
ступень наследует проверенные механизмы предыдущей и добавляет ровно
одно новое явление. Такой порядок гарантирует: если на нижней ступени
что-то сломано, это видно сразу; если верхняя ступень даёт неожиданность,
её источник локализуется по контрасту.

\begin{center}
\small
\begin{tabularx}{\textwidth}{@{}l l l >{\raggedright\arraybackslash}X@{}}
\toprule
Ступень & $n$ & Род & Константа и примечание \\
\midrule
Тор & $4$ & $1$ & тройка $(4,1,4\pi^2)$; четыре спиновые структуры;
  фаза $\pi/4$; вывод поправок из первых принципов \\
Кривая Клейна & $7$ & $3$ & $b_{Ch}(7)\approx0.7526$; фазы
  $\pi/2,\pi/3,\pi/7$; $B=b_2(K3)=22$; $j=-3375$ \\
Макбит & $9$ & $7$ & $b_{Ch}(9)\approx0.1172$; башня Гурвица \\
Hurwitz$_3$ & $11$ & $14$ & $b_{Ch}(11)\approx0.0823$; строка согласована
  по $\mathrm{PSL}(2,13)$ \\
\textbf{Ферма (стенд)} & $\mathbf{15}$ & $\mathbf{91}$ &
  $b_{Ch}(15)$ — радикал~\eqref{eq:b15}; фазовая решётка
  $\tfrac1{15}\ZZ[\zeta_{15}]$ \\
\textbf{Ферма (стенд)} & $\mathbf{30}$ & $\mathbf{406}$ &
  $b_{Ch}(30)$ — радикал~\eqref{eq:b30}; фазовая решётка
  $\tfrac1{30}\ZZ[\zeta_{30}]$ \\
\bottomrule
\end{tabularx}
\end{center}

Лестница демонстрирует последовательное наращивание масштаба: род растёт
от $1$ к $406$, фазовые решётки усложняются от $\ZZ[i]$ до
$\ZZ[\zeta_{30}]$, а целочисленных индексов становится больше — по блокам
проводников. При этом сами правила не меняются: спиновая фаза всегда
$\pi/n$, торможение всегда $\delta^4/k$, нормировка всегда строится из
самих периодов. Постоянство правил при росте масштаба — главный
методологический результат лестницы.

\begin{remark}[о роли вычислительных сертификатов]
Каждое количественное утверждение монографии снабжено сертификатом двух
возможных типов. Сертификат первого типа — точная целочисленная
арифметика: переписи, ранги, определители, индексы вычисляются без
какого-либо округления и проверяются независимым алгоритмом. Сертификат
второго типа — прямое независимое интегрирование: тождества с
$\Ga$-функциями сверяются квадратурированием tanh-sinh по гладким
кускам, в котором $\Ga$-функции не участвуют вовсе. Вычислительный
протокол лаборатории воспроизводит все сертификаты одной командой;
никаких «подсмотренных» констант в коде нет.
\end{remark}
